\begin{proposition}[The finite-field coefficient-duality criterion]
\label{prop:B-finite-field-coefficient-gate}
\status{Proved equivalence for all Chow-heart coefficients; the universal hypothesis remains conjectural.}
Fix a finite field \(k\), and let \(\mathcal H\) be the rational Chow
heart in \(\mathcal D=DM_{\rm gm}(k,\Q)\).  For a smooth projective
\(X/k\) of dimension \(d\), use the conventions
\[
 h_{\rm coh}(X)=M(X)^\vee=M(X)(-d)[-2d],\qquad
 P_{X,n}=h_{\rm coh}(X)(n)[2n].
\]
Thus \(\mathcal H\) is generated by the \(P_{X,n}\) under finite sums
and direct summands, and is closed under duality.  Use the prescribed
rational Weil extension of ordinary geometric motives, normalized by
a generator \(e\) of degree one at \(\Spec k\).  The following are
equivalent.
\begin{enumerate}
\item For every \(P\in\mathcal H\) and every integer \(i\), the
evaluation--trace pairing
\[
 W^i(P)\times W^{1-i}(P^\vee)\longrightarrow\Q
\]
is perfect between finite-dimensional spaces.
\item Beilinson's rational--numerical equality and Parshin's rational
vanishing hold for every smooth projective variety over \(k\).
\item For every \(P\in\mathcal H\), the prescribed constant coefficient
on \(\mathbf P^1_k\) has perfect finite-dimensional pairings
\[
 W^i(\mathbf P^1,\pi^*P)\times
 W^{3-i}(\mathbf P^1,\pi^*(P^\vee(1)))\longrightarrow\Q
 \quad\text{for every }i,
\]
with the usual projective-bundle formula, cup product and Gysin trace.
\end{enumerate}
Here ``Chow heart'' does not mean a single homogeneous cohomological
weight.  To infer the universal criterion from an assertion restricted
to homogeneous pure Chow summands, one must justify extension to the
whole Chow heart; an actual finite Chow--K\"unneth splitting supplies
such a justification.
\end{proposition}
\begin{proof}
Put \(A^i(P)=\Hom_{\mathcal D}(\mathbf1,P[i])\).  Our conventions give
\[
 P_{X,n}^\vee=P_{X,d-n},\qquad
 A^i(P_{X,n})=H_M^{2n+i}(X,\Q(n))
             =\operatorname{CH}^n(X,-i)_\Q.
\]
Since the higher-Chow complex has no negative simplicial degrees,
\(A^i(P)=0\) for \(i>0\), first on these generators and then on every
\(P\in\mathcal H\).  The same applies to \(P^\vee\).

The rational Weil splitting, together with its multiplicative
projection formula, gives a model
\[
 C_W(P)=E\otimes_\Q R\Hom_{\mathcal D}(\mathbf1,P),\qquad
 E=\Q\oplus\Q e,\quad |e|=1,\quad e^2=0,
\]
and hence
\[
 W^i(P)=A^i(P)\oplus e A^{i-1}(P).
\]
These are the rational coefficient formulas in Kahn's
Remark~10.5.2(2), Theorem~11.2.1 and Corollary~11.2.2(c)
\cite[pp.~47--49]{BKahn2012FullFaithfulness}.  For clarity, place
\(e\) on the left.  If \(a+eb\) and \(c+ef\) have total degrees
\(i\) and \(1-i\), the trace is exactly
\[
 b_P^i(a+eb,c+ef)=(-1)^i\operatorname{ev}(a,f)
                         +\operatorname{ev}(b,c).
\]
The sign comes from moving the degree-\(i\) element \(a\) past \(e\).
Both evaluations on the right have ordinary degree zero.

Suppose (1) holds.  For \(i<0\), negativity gives
\(W^{1-i}(P^\vee)=0\).  Perfection forces \(W^i(P)=0\), and therefore
its direct summand \(A^i(P)\) vanishes.  Together with negativity,
\(A^i(P)=0\) for every \(i\ne0\).  Taking \(P=P_{X,n}\) is precisely
Parshin's vanishing in every twist.  In degrees zero and one the
pairing is now ordinary evaluation
\[
 \operatorname{CH}^n(X)_\Q\times
 \operatorname{CH}^{d-n}(X)_\Q
 \longrightarrow\Q,\qquad
 (\alpha,\beta)\longmapsto\deg(\alpha\beta).
\]
Its kernel on either side is the numerically trivial subgroup, so
perfection gives Beilinson's equality.

Conversely assume (2).  Parshin and negativity make \(A^i\) vanish
outside degree zero on every generator.  Beilinson makes the displayed
intersection pairing nondegenerate.  Numerical cycle groups are
finite dimensional: any independent finite family of numerical classes
admits a full-rank rational intersection matrix, whose rank is bounded
by the corresponding finite Weil-cohomology dimension.  The pairing
is consequently perfect.  This persists under finite sums and direct
summands, since the dual idempotent is adjoint to the original
idempotent under evaluation.  Thus only
\(W^0(P)=A^0(P)\) and \(W^1(P)=eA^0(P)\) remain, and the two cross
pairings are perfect.  This proves (1).  It is the Chow-heart version
of Kahn's Proposition~10.5.1 and Remark~10.5.2(2), whose argument
also uses exactly the negative higher-Chow degrees and numerical
intersection pairing \cite[pp.~46--47]{BKahn2012FullFaithfulness}.

For the curve statement put \(h=c_1(\mathcal O(1))\), with degree two
and twist one.  The projective-bundle formula reads
\[
\begin{aligned}
 W^i(\mathbf P^1,\pi^*P)
   &=W^i(P)\oplus h W^{i-2}(P(-1)),\\
 W^{3-i}(\mathbf P^1,\pi^*(P^\vee(1)))
   &=W^{3-i}(P^\vee(1))\oplus hW^{1-i}(P^\vee).
\end{aligned}
\]
Since \(\pi_*h=1\) and \(h^2=0\), the trace pairing is block
off-diagonal.  On the two pairs of components it is
\[
 ((x,u),(y,v))\longmapsto
 b_P^i(x,v)+b_{P(-1)}^{i-2}(u,y).
\]
There is no extra sign from \(h\), whose degree is even.  The
unit--unit component is killed by the Gysin projection, and the
\(h\)--\(h\) component is zero.  Thus (3) implies (1) on the first
cross-block, including finite dimensionality.  Conversely
\(P'=P(-1)[-2]\) belongs to \(\mathcal H\), with
\[
 W^i(P')=W^{i-2}(P(-1)),\qquad
 W^{1-i}((P')^\vee)=W^{3-i}(P^\vee(1)).
\]
Condition (1) for \(P\) and \(P'\) proves perfection of both blocks,
and hence (3).

Finally, the argument uses the prescribed ordinary-motive Weil
extension.  It does not replace those motivic groups by arbitrary
abstract dual completions.  Nor does it assume that the image of
ordinary motives is itself a triangulated subcategory of the larger
Weil category; Kahn's Remark~11.3.2 explicitly explains this failure.
A cohomological K\"unneth projector does not by itself provide an
actual Chow--K\"unneth splitting.  This proves the stated limitation
for a formulation involving only homogeneous pure coefficients.
\end{proof}

Under B1's standing universal Tate and semisimplicity hypotheses, the
verified Geisser--Kahn implication supplies Parshin once universal
Beilinson holds.  In that setting the preceding criterion retains the
same Chow-phantom gate.  Without those standing hypotheses it is the
joint Beilinson--Parshin condition.  This finite-field dependency is
separate from the number-ring f-admissibility condition below; neither
follows from finite-coefficient duality or the proved Artin--Tate cases.
